\section{Other formal systems and concise encodings}
\label{sec:encodings}

The successful state reduction uses the K theory already calibrated by
Glazer. The broader search nevertheless produced useful results about
what to encode and how to charge its cost. A short mathematical formula,
a short complete data packet, a small register program and a small
blank-start binary machine are four different optimization objectives.
There is no fixed conversion factor between them.

\subsection{Comparison of the explored routes}

\begin{center}\small
\begin{tabularx}{\linewidth}{@{}p{30mm}XX@{}}\toprule
Route & Advantage & Current outcome and cost boundary\\\midrule
Direct ZFC proof search & The target equivalence follows directly by
enumerating finite derivations. & A baseline conceptual construction.
Schema recognition, substitutions, variable conditions and proof checking
all have to be implemented; no competitive small literal table was
constructed by this route.\\
Compact K schema & One closure-and-witness schema replaces a larger
set-theory packet, while retaining the numerical calculus. & The route
used by the selected machine. Both finite proof translations and the
actual emitter are available in the pinned Rocq source.\\
Finite equational noncollapse & A finite operation signature, finitely
many equations and one ground query replace axiom-schema recognition. &
E663 is inherited from the 235-state package. A complete new prefix packet
has 574 bits, including its codebook. No binary proof-search interpreter
or state bound for it is claimed.\\
Constructor fusion & One operation can serve several semantically
compatible syntactic roles. & The supplied E679 study gives an explicit
equiconsistent fused presentation, with complete expansion/interpretation
arguments. Its fixed-code packet is 679 bits; under the new prefix codec
it is 586 bits.\\
Finite class-theory frontend & A finite axiom list may simplify the
enumerator's axiom-recognition step. & An architectural alternative only.
The costs of class syntax, sorts, comprehension encoding and a complete
derivation translation were not converted into a competitive literal
machine in this investigation.\\
\bottomrule\end{tabularx}
\end{center}

The first route is mathematically straightforward: enumerate all finite
strings and halt precisely when one is a valid classical ZFC proof of
falsity. A computable bijective enumeration and a decidable proof checker
give the desired equivalence by elementary arithmetic. Its disadvantage
here is not logical strength but the size of the complete low-level
implementation. The K route already has a carefully optimized emitter
and interpreter, so semantic source transformations and state sharing
can immediately be measured in actual native states.

The equational route removes the need for an unbounded axiom schema at
the data level, but does not remove proof search. Its machine must
enumerate substitutions, contexts and applications of finitely many
equations, recognize the selected ground equality, implement all data
access and initialize its packet from a blank tape. Omitting that
interpreter or its blank writer would hide the main cost of a claimed
small machine.

\subsection{The inherited E663 packet}

The 235-state archive supplies E663: nine operation symbols, seven
equations and 150 equation-term nodes, with ground query
\[
                          Q(d)=d.
\]
The signature is
\[
 d:0,\quad N:2,\quad V:2,\quad P:2,\quad A:2,\quad Q:1,
 \quad W:2,\quad p_E:2,\quad p_\varepsilon:2.
\]
The meanings include a Boolean operation, environment access, term and
environment constructors, membership evaluation and a witness operation.
The actual equations, not just these descriptions, are preserved in
\path{formalisms/E663.json}; the complete original derivation and
calibration are in the unchanged 235-state archive~\cite{proveit235}.

The relevant assertion is nonderivability of the selected query:
\[
        \Con(\mathrm{ZFC})\quad\longleftrightarrow\quad
                  E663\nvdash_{\mathrm{eq}}Q(d)=d.
\]
This is a noncollapse problem. An ordinary equational theory always has
a one-element model, so saying that its axioms are simply ``consistent''
would not express the intended distinction. The query identifies the
collapse that corresponds to a refutation of the set-theoretic system.
An equational proof enumerator would halt on a derivation of that query.

In the inherited fixed-width codec, E663 costs 43 header bits,
608 axiom bits and 12 query bits, totaling 663. Its seven equations have
respectively
\[
                         14,6,9,12,12,17,80
\]
term nodes. The Boolean single-identity ingredient is inherited, rather
than a new Boolean-algebra axiom discovered here~\cite{mccune}.

\subsection{A fully declared 574-bit prefix packet}

The new codec allocates short codes to frequent term symbols, including
variables. The complete codebook is part of the packet, so the decoder
does not need a preinstalled symbol palette or an external arity list.
The packet contains thirteen codebook leaves: nine operations and four
variable slots. A binary prefix tree is serialized recursively using
the following self-delimiting node tags:
\begin{center}
\begin{tabular}{@{}ll@{}}\toprule
Tag & Meaning\\\midrule
\code{0} & Internal branch; two child codebooks follow\\
\code{10} & Leaf denoting a binary operation\\
\code{110} & Leaf denoting a variable slot\\
\code{1110} & Leaf denoting a constant\\
\code{1111} & Leaf denoting a unary operation\\
\bottomrule\end{tabular}
\end{center}
Operation and variable identities are their respective preorder indices
among the leaves. The path from the root to a leaf is its term-token
code. Once a token is read, its arity determines how many subterms follow.
Thus a term is decoded recursively without a terminator after each node.

An equation is a one bit followed by its two terms. A zero bit ends the
equation list and is followed by the two ground terms of the query.
The decoder requires the packet to end exactly there. Variable names
are local to an equation, as they are universally quantified there;
renaming them by occurrence frequency does not change equational
consequence. Operation names are uniformly renamed, with arities
preserved. There is no identification of distinct operations or variables.

For the selected E663 packet the complete codebook is:
\begin{center}\small
\begin{tabular}{@{}llrr@{}}\toprule
Symbol & Prefix code & Frequency & Arity/kind\\\midrule
$Q$ & \code{0000} & 7 & 1\\
$p_E$ & \code{0001} & 7 & 2\\
$V$ & \code{0010} & 8 & 2\\
$v_3$ & \code{00110} & 4 & variable\\
$W$ & \code{001110} & 2 & 2\\
$A$ & \code{001111} & 3 & 2\\
$v_0$ & \code{01} & 33 & variable\\
$v_1$ & \code{100} & 19 & variable\\
$N$ & \code{101} & 21 & 2\\
$p_\varepsilon$ & \code{1100} & 10 & 2\\
$v_2$ & \code{1101} & 12 & variable\\
$d$ & \code{1110} & 13 & 0\\
$P$ & \code{1111} & 14 & 2\\
\bottomrule\end{tabular}
\end{center}
The $v_i$ labels in this table are the encoder's frequency-ranked slots;
the raw decoder may name them in preorder instead. The included audit
checks precisely that bijective renaming for each equation.

The total charged cost is
\[
 \underbrace{46}_{\text{codebook}}
 +\underbrace{508}_{\text{equation terms}}
 +\underbrace{8}_{\text{equation framing}}
 +\underbrace{12}_{\text{ground query}}
 =\boxed{574\text{ bits}}.
\]
This is 89 bits shorter than the inherited 663-bit encoding. Byte
transport requires 72 bytes, with two padding bits that are not part of
the packet. The bit-character SHA--256, excluding the file's final
newline, is
\begin{quote}\small
\path{c6414e72e7682442c170ef20c3f0376ee5ee82368854ec231bc500c565ae6313}.
\end{quote}

\begin{proposition}[Exact presentation preservation]
Decoding the delivered 574-bit packet, using only that packet and the
fixed codec grammar, yields exactly E663 and its ground query up to a
bijective operation renaming and independent bijective variable
renaming in each equation. In particular, the decoded query is derivable
from the decoded equations if and only if the original query is derivable
from E663.
\end{proposition}
\begin{proof}
The codebook parser is deterministic because its five node tags are
prefix-free. Each leaf is assigned one identity and one arity. Induction
on a term shows that token decoding followed by decoding its arity-many
children reverses its prefix encoding. Induction on the equation list
and its final zero delimiter extends this to the complete packet and
the two ground terms. The retained round-trip calculation compares
every decoded term with the specified renaming of the original term.
Apply the operation bijection uniformly throughout a derivation. For
each axiom separately, its variable bijection preserves exactly the
set of substitution instances: compose a substitution with that
bijection or its inverse. The two presentations therefore have
corresponding axiom instances, and the unchanged rules for contexts,
symmetry, reflexivity and transitivity derive corresponding ground
equalities. The inverse operation renaming gives the converse.
\end{proof}

The independent decoder accepts only the bit string; it does not read
the encoder's JSON signature, token codes or variable maps. The replay
also rejects all 574 strict truncations, both one-bit extensions and
a nonbinary-character control. These diagnostics support the exact
round trip; the mathematical preservation argument is the induction
above.

\subsection{Why the shorter packet is not yet a state bound}

A Huffman codebook reduces data length but asks more of the runtime
parser than a fixed-width alphabet. A very small interpreter might
prefer a slightly longer packet whose tokens can be processed with
fewer states. A blank writer can also reuse repeated substrings rather
than paying a uniform state cost per bit. Consequently the 89 saved
bits cannot be subtracted from a native state count.

The concrete next experiment for this route is a complete binary
equational-proof enumerator with a counted decoder and blank writer.
Only its resulting full transition table would support comparison
with the selected K machine. The artifact records a null TM state
bound for the prefix packet so that its smaller data cost is not
mistaken for such a construction.

The supplied E679 constructor-fusion study illustrates the complementary
semantic route. It reduces the operation vocabulary by arranging that
one operation has compatible roles in a model with fresh copies of
each membership type. That argument concerns interpretations and
equiconsistency; the prefix codec concerns lossless representation of
an already fixed presentation. The methods can be combined, but they
must be charged and proved separately. Under the same new codec,
E679 occupies 586 bits, so E663's 574-bit packet is the better of the
two concrete packets considered here.
